\documentclass[10pt,a4paper]{amsart}
\usepackage[margin=19mm]{geometry}
\usepackage{amsmath,amssymb,mathtools}
\usepackage[scr=rsfs,cal=euler]{mathalfa}
\usepackage{xfrac}
\usepackage[unicode,hidelinks,bookmarks=false]{hyperref}
\newcommand{\ie}{i.e.\ }
\newcommand{\cf}{cf.\ }
\newcommand{\resp}{respectively\ }
\newcommand{\Subsection}{\subsection}
\providecommand{\nobreakdash}{\nobreak\hbox{-}}
\setlength{\emergencystretch}{2em}
\multlinegap0pt
\usepackage{amssymb}
\usepackage{enumerate}
\usepackage{xcolor}
\usepackage{tikz}
\newtheorem{thm}{Theorem}
\def\GL{{\rm GL}}
\title{Triality and Adjoint lifting for $\GL_3$}
\author{Wee Teck Gan}
\date{Source: arXiv:2512.08307v3 (CC BY 4.0)}
\begin{document}
\maketitle

\noindent\emph{Source and licence.} Triality and Adjoint lifting for $\GL_3$. Version arXiv:2512.08307v3 (CC BY 4.0), distributed by arXiv under the Creative Commons Attribution 4.0 International licence, http://creativecommons.org/licenses/by/4.0/, as recorded in the arXiv licence metadata for this exact version and shown on the abstract page https://arxiv.org/abs/2512.08307v3. Full licence notice: \texttt{licenses/arxiv-2512.08307-CC-BY-4.0.txt}.\par\smallskip

\noindent\textit{Editorial scope.} Counted unit: the article's thm T:cuspidality (source0.tex lines 2091--2103). The article's complete original proof of it is delivered with it (source0.tex lines 2245--2253, one \texttt{proof} environment of 9 lines, which the article places away from the statement). No mathematics is written, repaired or re-derived: the statement and the proof are the article's own lines, in the article's own notation, and no retained line is substituted or re-flowed.  No further source-local context is needed: every cross-reference of every retained line resolves inside the two delivered spans.  Nothing was omitted from any retained span, so the package carries no omission record.  No retained line contains a citation, so the wrapper carries no bibliography and no bibliography entry.  No retained line contains a figure environment, an image-inclusion command or drawing code, so no image is delivered and the package declares no asset dependency.  Editorial formatting only, in addition: an AMS \texttt{amsart} preamble that restates the article's own declarations and macros used by the retained lines, namely the environments \texttt{thm} on separate counters, together with the standard packages those lines need.  The retained environments print in this document as Theorem 1 in order of appearance, whereas the article prints the same items with the numbering of its own full document.  The complete native source of the article as distributed in the arXiv e-print for this exact version is bundled unchanged as \texttt{sources/190648/source0.tex}.
\begin{thm} \label{T:cuspidality}
Let $\pi$ be a unitary cuspidal automorphic representation of $\GL_3$ which is not an automorphic induction from a cubic field extension or a twist of a self-dual representation. 
Assume that the above hypothesis holds for $\pi$. Then ${\rm Ad}(\pi)$ is a self-dual automorphic representation of orthogonal type with trivial central character. Moreover,
 one of the following holds:
\vskip 5pt
\begin{itemize}
\item[(i)]  ${\rm Ad}(\pi)$ is  cuspidal  and $L^S(s, {\rm Ad}(\pi), \wedge^3)$ has a simple pole at $s=1$.
\vskip 5pt

\item[(ii)]  ${\rm Ad}(\pi) = A \boxplus B$ with $A$ and $B$ self-dual cuspidal representations of orthogonal type with the same  central character which is a nontrivial quadratic character. If $K/k$ is the quadratic field extension determined by the central character of $A$,  then
$\pi_K$ is dihedral with respect to four distinct non-Galois cubic extensions of $K$.  Moreover, $L^S(s, {\rm Ad}(\pi), \wedge^3)$ has a double pole at $s=1$.
\end{itemize}
\end{thm}

\begin{proof}
If  $\pi \boxtimes \tau_1 \simeq \pi \boxtimes \tau_2$, then the partial L-function
\[  L^S(s,  (\pi \boxtimes \tau_1) \times (\pi \boxtimes \tau_2)^{\vee}) = L^S(s, \pi \times \pi^{\vee} \times \tau_1 \times \tau_2^{\vee}) \]
has a pole at $s=1$.  But this L-function is equal to
\[  L^S({\rm Ad}(\pi) \times (\tau_1 \boxtimes \tau_2^{\vee})) \cdot  L^S(s, \tau_1 \times \tau_2^{\vee}). \]
Under the hypothesis on $\pi$, it follows by Theorem \ref{T:cuspidality} that ${\rm Ad}(\pi)$ is either cuspidal or an isobaric sum of two self-dual cuspidal representations of $\GL_4$ with nontrivial central character.
This implies that the L-function $L^S({\rm Ad}(\pi) \times (\tau_1 \boxtimes \tau_2^{\vee}))$ is holomorphic at $s=1$, since the self-dual representation $\tau_1 \boxtimes \tau_2$ has trivial central character. 
Hence, the L-function $L^S(s, \tau_1 \times \tau_2^{\vee})$ must have a pole at $s=1$, which implies that $\tau_1 \simeq \tau_2$.
\end{proof}

\end{document}
